\subsection{Compositional Pattern Recall (TinyStories)}
\label{sec:compositional}

\paragraph{Task Description.}
The compositional experiment evaluates whether the substrate can reconstruct
the ending of a short story based on structural patterns learned from other
stories.  The corpus is \texttt{roneneldan/TinyStories} (100 ingested samples):
a collection of English short stories for children, characterised by highly
regular grammar (''Once upon a time there was a [adj] [noun] who liked to
[verb]\ldots'') and controlled vocabulary with substantial cross-story overlap.
The held-out target (rightmost 30\% of each sample) must be reconstructed
purely from chord attractor resonance over the ingested story patterns.

\paragraph{Results.}
Across $N = 10$ test samples the mean weighted score was 0.000
(exact: 0.0\%, partial: 0.000).
\paragraph{Assessment.}
Recall quality was low.  At 100 ingested samples the substrate has not yet
accumulated sufficient TinyStories pattern density to reliably reconstruct
held-out story endings.  A larger ingestion corpus will yield clearer results.

Figure~\ref{fig:compositional_map} shows the trial outcome map.
